\documentclass{article}
\usepackage{geometry}
\usepackage{concmath}
\usepackage[OT1]{fontenc}
\usepackage{graphicx}
\usepackage{amsfonts}
\usepackage{amsmath,mathtools}
\usepackage{calrsfs}
\usepackage{color}
\usepackage[dvipsnames]{xcolor}
\usepackage[many]{tcolorbox}
\usepackage{setspace}
\usepackage{tkz-tab}               % for LINE SPACING
\usepackage{multicol}  
\geometry{margin=0.5in}
\definecolor{green}{RGB}{52, 235, 168}
\definecolor{violet}{RGB}{165, 83, 237}
\definecolor{lemon}{RGB}{170, 237, 83}
\definecolor{red}{RGB}{240, 125, 84}
\definecolor{blue}{RGB}{105, 255, 242}

\begin{document}
\noindent
\colorbox{green}{Ex 1}\\
\underbar{\(2 \Rightarrow 1\)}\\
Si \(u = w^k\) et \(v=w^p\) alors \(uv = w^p w^k\) = vu \\
\underline{\(1 \Rightarrow 2\)}\\
H(n) : "Si uv = vu et \(|u|+|v| = n\) alors \(\exists w, k, p, u = w^k, v = w^p\)"\\
H(0) est vraie car si \(|u| + |v| = 0\) alors \(u = v = \varepsilon\)
et \(w,k,p = \varepsilon, 1, 1\) convient\\
Soit \(n \in \mathbb{N}^*\)\\
Supposons H(k) pour \(k < n\)\\
Soit u,v tq uv = vu et \(|u| + |v| = n\)\\
* Cas  \(|u| > |v|\) idem cas symétrique\\
uv = vu \(\Rightarrow \exists u', u = vu'\)\\
vu'v = vvu'\\
\(\Rightarrow u'v = vu'\)\\
Si \(v = \varepsilon\) alors w,k,p = u,1,0 convient\\
Sinon, \(|u'| < |u|\)\\
donc \(|u'| + |v| < |u|\)
Donc \(|u'| + |v| < n\)\\
Par récurrence, \(\exists w,k,p\) tq \\
\(u'=w^k, v = w^p\)\\
et \(u = vu' = w^{p+k}\)\\
D'où H(n)\\
Si \(|u| = |v| : uv = vu \Rightarrow u = v\)\\
et w, k, p = u, 1, 1 convient 

\noindent \colorbox{green}{Ex 2}\\
1/\\
H(n) : "..."\\
H(2) et H(3) sont vraies car \(f_2 = ba\) et \(f_3 = bab\)\\
Supposons H(n) pour un \(n \geq 2\)\\
Mq H(n+2)\\
Alors \(f_{n+2} = f_{n+1}f_n\) a le même suffixe que \(f_n\) car
\(|f_n| \geq 2\)\\
2/\\
\underline{Initialisation}\\
H(3), H(4), H(5)\\
\underline{Hérédité}\\
Soit \(n \in \mathbb{N}\) tq H(n), H(n+1), H(n+2)\\
\begin{align*}
    f_{n+3} &= f_{n+2}f_{n+1}\\
    &= f_{n+1}f_n f_{n+1}
\end{align*}

\noindent \colorbox{green}{Ex 3}\\
1/ Faux avec e = a car \(a \in L(e*)^2\) (\(a = a \varepsilon\))
mais \(a \notin L(e^2)* = L(aa)*\)\\
2/ Faux : si \(e_1 = a, e_2 = b\) alors \begin{align*}
    ab &\in (a|b)*\\
    ab &\notin a*|b*
\end{align*}
3/ \((ab)* \neq a*b*\)\\
4/ \(\supseteq :\) Si \(u \in (e_1* e_2*)*\)
alors u est une concatanénation de mots de \(e_1\)
et \(e_2\) donc \(u \in (e_1|e_2)*\)
\(\subseteq :\) Si \(u \in (e_1 |e_2)*\) : 
en décomposant u suivant des blocs consécutifs de \(e_1\): \\
\(u \in e_1...e_1e_2...e_2e_1...e_1...\)\\
\(u \in e_1^{k_1}e_2^{k_2}...\)\\
\(u \in (e_1*e_2*)*\)

\noindent \colorbox{green}{Exercice 4}\\
1. \begin{equation*}
    (c*ac*bc*|c*bc*ac*)
\end{equation*}
2. \begin{equation*}
    (ab|ac|bc)*(a|\varepsilon)
\end{equation*}
3. \begin{equation*}
    (b|bc)*a(b|bc)*a(b|bc)*
\end{equation*}
4. \begin{eqnarray*}
    L(\frac{1}{6}) = 1(6)*|\varepsilon\\
    L(\frac{1}{7}) = (142857)*(1|14|142|1428|14285|\varepsilon)
\end{eqnarray*}
5. Si \(x \in \mathbb{Q} : x=x_1, x_2 pp...\) p: plus petite période\\
Soit Pref(p) l'ensemble des préfixes de p.\begin{align*}
    L(x)&= Pref(x_2)|x_2p*Pref(p)
\end{align*}
fini donc régulier
6. \begin{equation*}
    (a^2)*(b^2)*|(a^2)*a(b^2)*b
\end{equation*}

\noindent \colorbox{green}{Exercice 5}\\
1. \begin{itemize}
    \item d(u,v) = d(v,u)\\
    \item d(u,v) = 0 ssi u = v\\
    \item Mq d(u,w) = d(u,v) + d(v,w)\\
    \begin{align*}
        u_i \neq w_i \Rightarrow u_i \neq v_i \text{ ou } v_i \neq w_i\\
        \text{ens }
    \end{align*}
\end{itemize}
2. \begin{equation*}
    0*1*|0*10*1*|0*1*0*
\end{equation*}
3. P(L) : "H(L) est régulier"\\
* Si L est fini alors H(L) finis (nb finis de modif pour chaque mot de L) donc régulier\\
* Si \(P(L_1), P(L_2)\): \begin{align*}
    &H(L_1 \cup L_2) = H(L_1) \cup H(L_2) \text{ modifer un mot de \(L_1\cup L_2\) revient à modifier un mot de \(L_1\) ou \(L_2\)}\\
    &H(L_1\cup L_2) \text{ regulier}\\
    &H(L_1L_2) = L_1H(L_2) \cup H(L_1)L_2 \text{ regulier}\\
    &H(L_1*) = L_1*H(L_1)L_1* \cup \{\varepsilon\} \text{ regulier}
\end{align*}
Donc \(\forall L\) régulier, P(L)

\noindent \colorbox{green}{Exercice 6}\\
1. \begin{eqnarray*}
    \overline{L_0} = \Sigma* a \Sigma* a\Sigma*\\
    \overline{L_1} = \Sigma*
\end{eqnarray*}
2. \(\overline{L} \subset \overline{\overline{L}}\): Soit \(u \in \overline{L}\)\\
\(u \prec u \) donc \(u \in \overline{L}\)\\
\(\overline{\overline{L}} \subset \overline{L}\), Soit \(u \in \overline{\overline{L}}\)\\
\(\exists, \varphi tq \forall i, v_i = u_{\varphi(i)}\)\\
\(\exists w, \psi\) tq \(\forall i, w_i = v_{\psi(i)} = u_{\varphi \circ \psi (i)}\)\\
\noindent
3. L' un langage fini, \(\overline{L}\) est toujours infini\\
4. Si L fini \begin{equation*}
    \overline{L} = \bigcup_{u\in L, u = u_1,...,u_n} \Sigma*u_1\Sigma*...\Sigma*u_n\Sigma*
\end{equation*}
* \(P(L_1), P(L_2)\):
\begin{align*}
    \overline{L_1 \cup L_2} = \overline{L_1} \cup \overline{L_2} \text{ regulier}\\
    \overline{L_1 L_2} = \overline{L_1}\overline{L_2}\\
    \overline{L_1*} = \Sigma *
\end{align*}

\noindent \colorbox{green}{Exercice 1}\\
1/ Soit L reconnu par A = \((\Sigma, Q, I, F, E)\)\\
Alors \(\tilde{L} = L(A')\) où \(A'=(\Sigma, Q, F, I, E')\)
où \(E'=\{(q,a,q')|(q',a,q)\in E\}\)\\
car \(u \in L(A') \Leftrightarrow \exists\) chemin \(q_0 \to q_1 ... \to q_n\) dans A'\\
chemin dans A
\(u_n...u_1 \in L\)
donc \(L(A') = \tilde{L}\)
2/ \(a*(b|\varepsilon)\)\\
\((ab)*(a|\varepsilon)\)
3/ Soit A = \((\Sigma, Q, I, Q', \delta)\) où Q' = l'ensemble des 
états co-accessibles\\
A''=(\(\Sigma, Q, Q''(\text{ états accessibles }), F, \delta\)) : reco Suff(L)
A'''=\((\Sigma, Q, Q'', Q', \delta)\) reco Fact(L)
4/ H(L)"Pref(L) régulier"\\
Si L fini : Pref(L) est fini, donc P(L)
Si Pref(\(L_1\)) \(Pref(L_2)\) réguliers\\
\(Pref(L_1 \cup L_2) = Pref(L_1) \cup Pref(L_2)\) régulier\\
\(Pref(L_1L_2) = Pref(L_1)\cup L_1 Pref(L_2)\) régulier\\
\(Pref(L_1^*) = L_1^*Pref(L_1)\) régulier\\
\(Fact(L_1 L_2) = Fact(L_1) \cup Fact(L_2) \cup Suff(L_1)Pref(L_2)\)\\
\(Fact(L*) = Suff(L)L^*Pref(L) \cup Fact(L)\)

\end{document}